\section{Summary}\label{sec:offline-summary}

This chapter addresses offline retail CBDC payments. These are payments
conducted while both payer and payee are disconnected from the network, and
settled only when tokens are eventually deposited. The central difficulty is
that a payee has no way to verify, in real time, that the token they are
receiving has not already been spent elsewhere. Rather than rely on prevention
or detection alone, we combine the two. A secure element makes double spending
expensive and impractical to mount at scale. An efficient auditing mechanism
ensures that, should an attack succeed, the responsible party can still be
identified and held accountable. Because the auditor names that party, the
double spent value is settled against them rather than absorbed by the issuer or
borne by whoever presented the second deposit, so no honest holder loses value
and the total in the system is unchanged.

The scheme keeps the institutions that an offline payment system already has,
and constrains what each of them learns. The central bank issues and redeems
tokens without learning what it signs at withdrawal or which withdrawal a
deposit corresponds to. Intermediaries learn the identity of a user only when
that user withdraws or deposits. The registration authority enrols users once
and plays no part in payments. The auditor recovers an identity only from a
payment that double spends a token, and the role can be distributed via
threshold decryption. None of these participants is removed in pursuit of
decentralisation, and none of them is trusted with the privacy of honest users.

We claim no new primitive. The contribution is a combination of anonymous
credentials, blind signatures, verifiable encryption, and a joint proof of
knowledge. The joint proof splits a payment between a secure element holding the
user's key and a mobile device doing the heavy work. What the combination buys
is that every property an offline payment system needs holds at once, on
hardware that exists today. These are unforgeability, theft prevention, double
spending identification, conservation of value, non-frameability, anonymity, and
unlinkability.

We demonstrate the viability of the protocol on commodity smart cards and mobile
phones. Withdrawal and the payer side of a payment both complete in under half a
second. The smart card itself authorises a payment in under $400\,\mathrm{ms}$.
Even a payment chain of 32 hops verifies in under a second on the payee's side.
Auditing adds no overhead for the payer, and a modest 10 to 20\% overhead for
the payee.

The credential infrastructure underlying this chapter is the same one used
in~\Cref{chap:txid}: the same registration, the same BBS+ credential, and the
same joint proof. A deployment can therefore identify transactions within an
epoch online, and identify double spenders offline, without maintaining two
separate sets of credentials. We give no joint security analysis of the two
running together, and a composable treatment of that combination is future work.

We focus on transferability rather than divisibility. Extending the scheme to
support splitting and merging tokens, and to support aggregated deposits, are
both natural directions for future work. Revocation in the offline setting
remains the sharpest limitation, and we discuss the expiry based mitigations we
consider practical in~\Cref{sec:offline-construction}.
